\chapter{How to use the templates} \label{ch-template}

This is a practical guide into how to use this template, by explaining the role of the different folders, and the options of \verb|\documentclass{oist_thesis}|, which accepts either \verb|temporary| or \verb|final|, plus an optional \verb|showframe|.

\section{Folders}

The repository contains two top-level folders. \textbf{Template} is this
document: the empty starting point you write your thesis in, and the one you
should copy. \textbf{Examples} is a complete thesis, kept so that you can look
up how something was actually done and copy the code across. The two share the
same \texttt{oist\_thesis.cls}, and \texttt{generate\_examples.sh} copies it
from \texttt{Template} into \texttt{Examples} so the formatting can never drift
apart.

Inside \texttt{Template} there are three folders, detailed here:

\begin{itemize}

	\item \textbf{Images.} This folder should contain all the images that you will use in your thesis. It can contain subfolders, for example one for each chapter. To include an image from the main text, use something like \texttt{\textbackslash includegraphics\{subfolder/image.jpg\} } without worrying about the \texttt{Images} path.

	\item \textbf{MainText.} This folder contains a series of \LaTeX\ files that form the main text: introduction, chapters, conclusion, appendices and published articles. The introduction and conclusion as they are now are not numbered, which creates a few difficulties with the headers of the thesis. Those are solved by including the commands \texttt{\textbackslash unnumberedchapter\{\}} and \texttt{\textbackslash numberedchapter} before including the files in \texttt{Thesis.tex}. If you want the introduction and conclusion to be numbered, include them as regular chapters.

	\item \textbf{Preamble.} This folder contains a series of \LaTeX\ files with the pages that will appear before the main text. Please write (or copy and paste) your own text in those files and delete the dummy text when appropriate. The files are:
	      \begin{itemize}
		      \item \texttt{abbreviations.tex} --- List of abbreviations. If the list goes over one page, create another table.
		      \item \texttt{abstract.tex} --- Abstract. Follow directions in the file.
		      \item \texttt{acknowledgments.tex} --- Acknowledgments. Follow directions in the file.
		      \item \texttt{declaration.tex} --- Declaration of Original and Sole Authorship. Insert the text from the template made available by the Graduate School in the thesis submission form.
		      \item \texttt{dedication.tex} --- Dedication (optional). Should only be a very few lines.
		      \item \texttt{glossary.tex} --- Glossary (optional). If the list goes over one page, create another table.
		      \item \texttt{nomenclature.tex} --- Nomenclature (optional). If the list goes over one page, create another table.

		            These last three are built with \texttt{xltabular} and a column
		            declared \texttt{X}, which is what makes the description column
		            stretch all the way to the right-hand margin before wrapping.
		            Add rows freely, but leave the table header line alone: replacing
		            \texttt{X} with a fixed width such as \verb|p{0.58\textwidth}|
		            makes the text wrap early and stop well short of the margin.
		      \item \texttt{physics\_bibstyle.bst} --- Bibliography style file modified by Jeremie Gillet in 2011 to suit his thesis. Might be suitable for physics. If you want to use another custom bibliography style, include the file in this folder.
		      \item \texttt{Thesis\_bibliography.bib} --- BibTeX file containing your bibliography.
		      \item \texttt{Thesis\_publications.bib} --- BibTeX file containing the list of your own publications (only for thesis by publication).
	      \end{itemize}

\end{itemize}

\section{\texttt{Thesis.tex}}

This is the main file, the only one that needs to be compiled to build the thesis. Compile once with \LaTeX, once with BibTeX and finally twice with \LaTeX\ to get all the references right. At the top of this file, you can see \verb|\documentclass[temporary]{oist_thesis}|. When you submit a temporary version to the graduate school, do not modify it. When you submit a final version, use \verb|\documentclass[final]{oist_thesis}| instead.

Let's go through each section and comment them briefly. The last section will emphasize the differences between options \verb|\documentclass[temporary]{oist_thesis}| and \verb|\documentclass[final]{oist_thesis}|.

\subsection{REQUIRED PACKAGES AND CONFIGURATIONS}

This section contains the minimum number of packages and definitions to compile the thesis. No line should be removed or modified.

\subsection{ADD YOUR PACKAGES}

Add you packages in this section.
In general, try to include as little packages as possible, as some of them can interfere and break each other.

\subsection{ADD YOUR DEFINITIONS AND COMMANDS}

Optionally, add custom definitions in this section.
You don't have to keep the examples that come with the template.

\subsection{PICK YOUR BIBLIOGRAPHY STYLE}

Pick the bibliography style the most adapted to your field within the \texttt{natbib} package.

\subsection{TITLE PAGE}

Creates the title page. Do not modify.

\subsection{PREAMBLE PAGES}

Structures the style (header) for the preamble pages and builds them. Do not modify, except for deleting the optional preambles you might not want to include.

\subsection{LIST OF CONTENTS/FIGURES/TABLES}

Creates the different lists. Do not modify.

\subsection{THESIS MAIN TEXT}

Structures the style for the main text chapters and builds them.

The command \texttt{\textbackslash numberedchapter} is only relevant for a transition between unnumbered sections and numbered sections, it does not need to be included between each chapter.

\subsection{OVERFULL LINES}

\LaTeX\ should warn you at compile time of overfull lines, with messages along the lines of
\begin{verbatim}
Overfull \hbox (14.45pt too wide) in paragraph at lines 5--151
\end{verbatim}
You should always address those warnings and fix the issues, usually by adding a \texttt{\textbackslash -} inside a word to indicate that it can be split at that point.
To make the issue more obvious, the template will add a thick black line in the margin wherever this happens, you might see several in this document.

In general, the final document should compile without any warnings or errors.

\subsection{OVER-WIDE EQUATIONS} \label{sec-wide-equations}

The black bar described above is drawn by \TeX\ itself, and it only ever appears for \emph{lines of text}.
Displayed mathematics is treated differently, in two ways that are worth knowing about, because a formula that is too long will happily run into the 2.5\,cm margins.

First, \LaTeX\ never breaks a display across lines on your behalf.
A paragraph of text is broken into lines automatically, but a formula is set as one long line, and if it does not fit it simply sticks out.
You have to choose the break point yourself.

Second, the black bar is unreliable for maths:
\begin{itemize}
	\item In \texttt{align}, \texttt{gather}, \texttt{multline}, \texttt{split}, \texttt{eqnarray} and in tables, \textbf{no bar is ever drawn}. These are built with \TeX's alignment machinery, which deliberately switches the bar off while it assembles them. No package or class can turn it back on.
	\item In \texttt{equation}, \texttt{equation*} and \verb|\[...\]|, a bar \emph{is} drawn, but at the end of the formula rather than at the edge of the text block. A formula that overshoots by more than the 2.5\,cm margin therefore draws its bar off the edge of the paper, where you cannot see it.
\end{itemize}

Whichever environment you used, the problem is always recorded in \texttt{Thesis.log}, so that is the place to check.
An equation reports as \texttt{Overfull \textbackslash hbox (62.31pt too wide) detected at line 88} and an \texttt{align} reports as \texttt{Overfull \textbackslash hbox (103.9pt too wide) in alignment at lines 120--126}.
Running \verb|./check_overfull.sh| after compiling collects all of these in one list.

Two things make over-wide maths easy to spot while you write.
Adding the \texttt{showframe} option, as in
\begin{verbatim}
\documentclass[temporary,showframe]{oist_thesis}
\end{verbatim}
draws the outline of the text block on every page, so anything sticking out is obvious; remove the option before you submit.
Setting \verb|\overfullrule=0pt| in \texttt{Thesis.tex} turns the black bars off again if you find them distracting, but do not submit a thesis that still needs them.

To fix a formula that is too long, break it yourself.
Use \texttt{split} inside an \texttt{equation} to spread one equation over several lines with a single number, breaking before a relation or a binary operator:
\begin{verbatim}
\begin{equation}
    \begin{split}
        f(x) &= a_1 x + a_2 x^2 + a_3 x^3 \\
             &\quad + a_4 x^4 + a_5 x^5 .
    \end{split}
\end{equation}
\end{verbatim}
Use \texttt{multline} when there is no natural alignment point, which sets the first line flush left and the last flush right, and \texttt{align} when you genuinely have several equations to line up.
Long fractions and long arguments to \verb|\left(...\right)| are the usual culprits, so it often helps to name a recurring subexpression with \verb|\newcommand| and use that instead.

Do not fix a wide equation with \verb|\resizebox|, \verb|\scalebox| or a smaller font.
Those change the size of the mathematics relative to the surrounding text, which the formatting requirements do not allow.

\subsection{OTHER THINGS THE TEMPLATE HANDLES FOR YOU}

A few more things are set up so that nothing drifts outside the 2.5\,cm margins.
You do not need to do anything about them, but it is worth knowing they are
there so you do not undo them by accident.

\begin{itemize}

	\item \textbf{Headings are ragged right.} A heading has very few places where
	      it can be broken, because the words are large, so when headings are
	      justified \TeX\ often prefers to overshoot the right margin rather than
	      open up an ugly gap. Ragged right simply breaks at the last word that
	      fits, so a long chapter title can never run into the margin.

	\item \textbf{URLs use the document font.} \texttt{hyperref} would otherwise
	      set them in typewriter, which looks out of place in the Reference List.
	      Because a proportional font gives fewer sensible break points, the
	      \texttt{xurl} package is loaded as well, so a long URL breaks wherever
	      it needs to instead of sticking out.

	\item \textbf{Numbered equations line up.} An equation number is placed
	      against \verb|\linewidth|, so an equation written inside anything that
	      narrows the text --- a \texttt{minipage}, a \texttt{quote}, a box of a
	      fixed fraction of \verb|\textwidth| --- gets its number pulled in from
	      the margin while every other equation sits flush with it. If your
	      equation numbers do not line up, look for a narrowing environment
	      around them rather than at the equation itself.

\end{itemize}

\subsection{BIBLIOGRAPHY}

Includes the bibliography.

\subsection{APPENDICES}

Structures the style for the appendices and builds them. The appendices are numbered with letters but are structured like regular chapters.

\subsection{Differences between temporary version and final version}

There is only one difference between
\begin{verbatim}
\documentclass[temporary]{oist_thesis}
\documentclass[final]{oist_thesis}
\end{verbatim}
namely the line spacing.
This might create small issues, such as page breaking with large tables, images or captions, when compiling the same content.

\subsection{The \texttt{showframe} option}

Besides \texttt{temporary} and \texttt{final}, the class accepts an optional \texttt{showframe}:
\begin{verbatim}
\documentclass[temporary,showframe]{oist_thesis}
\end{verbatim}
It draws the outline of the text block on every page so that you can see immediately whether an equation, a table or a figure spills into the margins, as discussed in Section~\ref{sec-wide-equations}.
It is a drafting aid only and changes no dimension of the document, but remove it before you submit.
